% *======================================================================*
%  Cactus Thorn template for ThornGuide documentation
%  Author: Ian Kelley
%  Date: Sun Jun 02, 2002
%
%  Thorn documentation in the latex file doc/documentation.tex
%  will be included in ThornGuides built with the Cactus make system.
%  The scripts employed by the make system automatically include
%  pages about variables, parameters and scheduling parsed from the
%  relevant thorn CCL files.
%
%  This template contains guidelines which help to assure that your
%  documentation will be correctly added to ThornGuides. More
%  information is available in the Cactus UsersGuide.
%
%  Guidelines:
%   - Do not change anything before the line
%       % START CACTUS THORNGUIDE",
%     except for filling in the title, author, date, etc. fields.
%        - Each of these fields should only be on ONE line.
%        - Author names should be separated with a \\ or a comma.
%   - You can define your own macros, but they must appear after
%     the START CACTUS THORNGUIDE line, and must not redefine standard
%     latex commands.
%   - To avoid name clashes with other thorns, 'labels', 'citations',
%     'references', and 'image' names should conform to the following
%     convention:
%       ARRANGEMENT_THORN_LABEL
%     For example, an image wave.eps in the arrangement CactusWave and
%     thorn WaveToyC should be renamed to CactusWave_WaveToyC_wave.eps
%   - Graphics should only be included using the graphicx package.
%     More specifically, with the "\includegraphics" command.  Do
%     not specify any graphic file extensions in your .tex file. This
%     will allow us to create a PDF version of the ThornGuide
%     via pdflatex.
%   - References should be included with the latex "\bibitem" command.
%   - Use \begin{abstract}...\end{abstract} instead of \abstract{...}
%   - Do not use \appendix, instead include any appendices you need as
%     standard sections.
%   - For the benefit of our Perl scripts, and for future extensions,
%     please use simple latex.
%
% *======================================================================*
%
% Example of including a graphic image:
%    \begin{figure}[ht]
% 	\begin{center}
%    	   \includegraphics[width=6cm]{/home/runner/work/tests/tests/arrangements/EinsteinInitialData/TwoPunctures/doc/MyArrangement_MyThorn_MyFigure}
% 	\end{center}
% 	\caption{Illustration of this and that}
% 	\label{MyArrangement_MyThorn_MyLabel}
%    \end{figure}
%
% Example of using a label:
%   \label{MyArrangement_MyThorn_MyLabel}
%
% Example of a citation:
%    \cite{MyArrangement_MyThorn_Author99}
%
% Example of including a reference
%   \bibitem{MyArrangement_MyThorn_Author99}
%   {J. Author, {\em The Title of the Book, Journal, or periodical}, 1 (1999),
%   1--16. {\tt http://www.nowhere.com/}}
%
% *======================================================================*

\documentclass{article}

% Use the Cactus ThornGuide style file
% (Automatically used from Cactus distribution, if you have a
%  thorn without the Cactus Flesh download this from the Cactus
%  homepage at www.cactuscode.org)
\usepackage{../../../../../doc/latex/cactus}

\newlength{\tableWidth} \newlength{\maxVarWidth} \newlength{\paraWidth} \newlength{\descWidth} \begin{document}

% The author of the documentation
\author{Marcus Ansorg \textless marcus.ansorg@aei.mpg.de\textgreater \\ Erik Schnetter \textless schnetter@cct.lsu.edu\textgreater}

% The title of the document (not necessarily the name of the Thorn)
\title{TwoPunctures}

% the date your document was last changed:
\date{May 12 2026}

\maketitle

% Do not delete next line
% START CACTUS THORNGUIDE

\begin{abstract}
TwoPunctures generates binary black hole initial data using the puncture
method with the Bowen-York extrinsic curvature ansatz.  The Hamiltonian
constraint is solved spectrally on a single computational domain that covers
the two asymptotically flat ends as a compactified radial coordinate, an
angular coordinate, and an azimuthal coordinate.  The resulting 3D metric,
extrinsic curvature, and lapse are stored in the ADMBase grid functions and
are ready for use by any evolution thorn that reads from ADMBase.  This thorn
targets the traditional Carpet AMR infrastructure.  A companion thorn,
TwoPuncturesX, provides the same functionality for the CarpetX
(AMReX-based) infrastructure.
\end{abstract}

\section{Introduction}

TwoPunctures computes conformally flat initial data for two non-spinning or
spinning, non-moving or moving black holes using the \emph{moving puncture}
approach \cite{einsteininitialdata_TwoPunctures_Brandt97}.  The central task
is to solve the Hamiltonian constraint for the conformal factor $\psi$, which
is decomposed as
\begin{equation}
  \psi = 1 + \frac{m_+}{2 r_+} + \frac{m_-}{2 r_-} + u,
\end{equation}
where $m_\pm$ are the bare masses of the two punctures located at
$(\pm b, 0, 0)$ (or offset by \texttt{center\_offset}), $r_\pm$ are the
coordinate distances from each puncture, and $u$ is a smooth correction
function that satisfies a nonlinear elliptic equation derived from the full
Hamiltonian constraint.  The Bowen-York extrinsic curvature
\cite{einsteininitialdata_TwoPunctures_BY80} is used to incorporate linear
momenta and spins.

The nonlinear elliptic equation for $u$ is solved with a spectral method
described in Ansorg, Br\"ugmann \& Tichy (2004)
\cite{einsteininitialdata_TwoPunctures_Ansorg04}.  Newton iteration is used
to handle the nonlinear terms.  Once $u$ is known, all ADMBase metric and
extrinsic-curvature components are filled on the 3D Cactus grid by
interpolation from the spectral solution.

\section{Physical System}

\subsection{Conformal Decomposition}

The spatial metric is assumed to be conformally flat:
\begin{equation}
  \gamma_{ij} = \psi^4 f_{ij},
\end{equation}
where $f_{ij}$ is the flat metric.  The extrinsic curvature is decomposed into
a trace-free part $A_{ij}$ and trace $K$:
\begin{equation}
  K_{ij} = \psi^{-2}\,A_{ij} + \tfrac{1}{3}\,\gamma_{ij}\,K.
\end{equation}
For the time-symmetric slice ($A_{ij}=0$ and $K=0$) the Hamiltonian constraint reduces to a
flat Laplace equation for $\psi$.  For the general Bowen-York case ($K=0$ is
still enforced on the slice but $A_{ij}\neq 0$), the Hamiltonian constraint
becomes a nonlinear equation for $u$.

\subsection{Bowen-York Extrinsic Curvature}

The Bowen-York extrinsic curvature for a single puncture with linear momentum
$P^i$ and spin $S^i$ is
\begin{equation}
  A_{ij}^{\rm BY} = \frac{3}{2r^2}\left[
    P_i n_j + P_j n_i - (f_{ij} - n_i n_j) P^k n_k
    + \frac{2}{r}\left(\epsilon_{kli} S^k n^l n_j
      + \epsilon_{klj} S^k n^l n_i\right)
  \right],
\end{equation}
where $n^i$ is the unit outward normal to a sphere of radius $r$ centered on
the puncture.  The total extrinsic curvature is the sum of the contributions
from both punctures.

\subsection{Masses}

The thorn distinguishes between \emph{bare masses} $m_\pm$ (the parameters
$\texttt{par\_m\_plus}$ and $\texttt{par\_m\_minus}$) and \emph{ADM masses}
$M_\pm$ measured at the respective asymptotic infinities.  When
\texttt{give\_bare\_mass = yes} the user sets bare masses directly.  When
\texttt{give\_bare\_mass = no} the thorn iterates to find the bare masses
that reproduce the requested target ADM masses \texttt{target\_M\_plus} and
\texttt{target\_M\_minus} to within the tolerance \texttt{adm\_tol}.

\section{Numerical Implementation}

\subsection{Spectral Grid}

The equation for $u$ is solved on a spectral grid parameterized by
coordinates $(A, B, \phi)$, where $A \in [0,1]$ is a compactified radial
coordinate covering both punctures and spatial infinity, $B \in [-1,1]$ is an
angular coordinate, and $\phi \in [0, 2\pi)$ is the azimuthal angle.  The
number of spectral coefficients in each direction is set by \texttt{npoints\_A},
\texttt{npoints\_B}, and \texttt{npoints\_phi}, respectively.

\subsection{Newton Solver}

Newton's method is applied to the nonlinear equation for $u$ with tolerance
\texttt{Newton\_tol} and at most \texttt{Newton\_maxit} iterations.  The
linear system at each Newton step is solved by a banded direct solver.
OpenMP parallelism is used within the Newton solver to accelerate the
construction and application of the Jacobian.

\subsection{Grid Interpolation}

Once the spectral coefficients are known, the solution is transferred to the
Cactus 3D grid in one of two ways controlled by \texttt{grid\_setup\_method}:
\begin{itemize}
  \item \textbf{Taylor expansion} (default): the solution is Taylor-expanded
    about the nearest spectral collocation point.  This is fast but may be
    slightly less accurate far from collocation points.
  \item \textbf{evaluation}: each grid point is evaluated using all spectral
    coefficients.  More accurate but significantly slower for large grids.
\end{itemize}
The interpolation loop is parallelized with OpenMP.

\subsection{Singularity Regularization}

Two small parameters help avoid numerical issues near the punctures:
\begin{itemize}
  \item \texttt{TP\_epsilon}: replaces $r_\pm$ with
    $(r_\pm^4 + \epsilon^4)^{1/4}$, smoothing the singularity.
  \item \texttt{TP\_Tiny}: a floor applied to $r_\pm$ after the above
    regularization.
  \item \texttt{TP\_Extend\_Radius}: inside this radius around each puncture
    the singular $1/r$ term is replaced by a smooth polynomial extension,
    useful when the puncture lies inside the computational domain.
\end{itemize}

\section{Using This Thorn}

\subsection{Obtaining This Thorn}

TwoPunctures is part of the \texttt{EinsteinInitialData} arrangement, which
can be checked out via the Einstein Toolkit release mechanism or directly from
the repository.  The thorn requires ADMBase, StaticConformal, and the
Cactus flesh.

\subsection{Basic Usage}

Add \texttt{EinsteinInitialData/TwoPunctures} to your thorn list and set the
following parameters in your parameter file:

\begin{verbatim}
ADMBase::initial_data  = "twopunctures"
ADMBase::initial_lapse = "twopunctures-antisymmetric"

TwoPunctures::par_b       = 3.0   # half separation (M)
TwoPunctures::par_m_plus  = 0.5
TwoPunctures::par_m_minus = 0.5
TwoPunctures::par_P_plus[1]  =  0.084
TwoPunctures::par_P_minus[1] = -0.084
TwoPunctures::npoints_A   = 30
TwoPunctures::npoints_B   = 30
TwoPunctures::npoints_phi = 16
\end{verbatim}

The punctures are placed at $(+b, 0, 0)$ and $(-b, 0, 0)$ unless
\texttt{center\_offset} shifts the center of mass.  To separate the holes
along the $z$-axis instead, set \texttt{swap\_xz = yes}.

\subsection{Lapse Options}

The initial lapse is selected via \texttt{ADMBase::initial\_lapse}:
\begin{itemize}
  \item \textbf{twopunctures-antisymmetric}: $\alpha =
    (1 - \psi_{\rm sing}/\psi) / (1 + \psi_{\rm sing}/\psi)$, ranges in
    $[-1, +1]$.
  \item \textbf{twopunctures-averaged}: average of the antisymmetric lapse
    and unity, ranges in $[0, +1]$.
  \item \textbf{psi\^{}n}: $\alpha = \psi^n$ where $n$ is set by
    \texttt{initial\_lapse\_psi\_exponent} (typically $-2$).
  \item \textbf{W}: $\alpha = \psi^{-2}$, suitable as a slow-start lapse for
    the W-formulation \cite{einsteininitialdata_TwoPunctures_Alic24}.
  \item \textbf{brownsville}: $\alpha = 2/(1 + \psi^n)$
    \cite{einsteininitialdata_TwoPunctures_Campanelli06}.
\end{itemize}

\subsection{Conformal Storage}

When \texttt{ADMBase::metric\_type = "static conformal"}, the thorn populates
the StaticConformal grid functions (\texttt{psi}, \texttt{psix}, etc.) in
addition to the physical metric, with the level of detail controlled by
\texttt{StaticConformal::conformal\_storage}.

\subsection{ADM Mass Iteration}

Setting \texttt{give\_bare\_mass = no} enables an outer iteration that finds
bare masses $m_\pm$ such that the resulting ADM masses match
\texttt{target\_M\_plus} and \texttt{target\_M\_minus} to within
\texttt{adm\_tol}.  The iteration uses the analytic inversion of the ADM
mass relation.

\subsection{Interaction With Other Thorns}

TwoPunctures writes to the ADMBase grid functions \texttt{gxx}, \texttt{gyy},
\texttt{gzz}, \texttt{gxy}, \texttt{gxz}, \texttt{gyz}, \texttt{kxx},
\texttt{kyy}, \texttt{kzz}, \texttt{kxy}, \texttt{kxz}, \texttt{kyz}, and
\texttt{alp}.  It exposes the scalar grid functions \texttt{mp}, \texttt{mm},
\texttt{mp\_adm}, \texttt{mm\_adm}, \texttt{E}, \texttt{J1}, \texttt{J2},
and \texttt{J3} for use by other thorns (e.g., Multipole, QuasiLocalMeasures).

Matter source terms can be supplied via the optional function interfaces
\texttt{Set\_Rho\_ADM} and \texttt{Set\_Initial\_Guess\_for\_u}, enabling
neutron star or BH-NS initial data when \texttt{use\_sources = yes}.

\section{History}

The spectral algorithm implemented in this thorn was developed by Marcus
Ansorg and is described in \cite{einsteininitialdata_TwoPunctures_Ansorg04}.
The thorn was originally written for the Cactus framework by Marcus Ansorg
and Erik Schnetter.  It has since been maintained and extended by the Einstein
Toolkit community.

\begin{thebibliography}{9}

\bibitem{einsteininitialdata_TwoPunctures_BY80}
  J.~M.~Bowen and J.~W.~York,
  {\em Time-asymmetric initial data for black holes and black-hole collisions},
  Phys.\ Rev.\ D {\bf 21}, 2047 (1980).

\bibitem{einsteininitialdata_TwoPunctures_Brandt97}
  S.~Brandt and B.~Br\"ugmann,
  {\em A simple construction of initial data for multiple black holes},
  Phys.\ Rev.\ Lett.\ {\bf 78}, 3606 (1997),
  gr-qc/9703066.

\bibitem{einsteininitialdata_TwoPunctures_Ansorg04}
  M.~Ansorg, B.~Br\"ugmann, and W.~Tichy,
  {\em Single-domain spectral method for black hole puncture data},
  Phys.\ Rev.\ D {\bf 70}, 064011 (2004),
  gr-qc/0404056.

\bibitem{einsteininitialdata_TwoPunctures_Campanelli06}
  M.~Campanelli, C.~O.~Lousto, and Y.~Zlochower,
  {\em Spinning-black-hole binary evolutions, recoil velocities, and
  spin-orbit effects},
  Phys.\ Rev.\ D {\bf 74}, 041501 (2006),
  gr-qc/0604012.

\bibitem{einsteininitialdata_TwoPunctures_Alic24}
  D.~Alic et al.,
  {\em Constraint-preserving boundary treatment for a harmonic formulation of
  the Einstein equations},
  Phys.\ Rev.\ D {\bf 110}, 064045 (2024),
  arXiv:2404.01137.

\end{thebibliography}

% Do not delete next line
% END CACTUS THORNGUIDE



\section{Parameters} 


\parskip = 0pt

\setlength{\tableWidth}{160mm}

\setlength{\paraWidth}{\tableWidth}
\setlength{\descWidth}{\tableWidth}
\settowidth{\maxVarWidth}{schedule\_in\_admbase\_initialdata}

\addtolength{\paraWidth}{-\maxVarWidth}
\addtolength{\paraWidth}{-\columnsep}
\addtolength{\paraWidth}{-\columnsep}
\addtolength{\paraWidth}{-\columnsep}

\addtolength{\descWidth}{-\columnsep}
\addtolength{\descWidth}{-\columnsep}
\addtolength{\descWidth}{-\columnsep}
\noindent \begin{tabular*}{\tableWidth}{|c|l@{\extracolsep{\fill}}r|}
\hline
\multicolumn{1}{|p{\maxVarWidth}}{adm\_tol} & {\bf Scope:} restricted & REAL \\\hline
\multicolumn{3}{|p{\descWidth}|}{{\bf Description:}   {\em Tolerance of ADM masses when give\_bare\_mass=no}} \\
\hline{\bf Range} & &  {\bf Default:} 1.0e-10 \\\multicolumn{1}{|p{\maxVarWidth}|}{\centering (0:*)} & \multicolumn{2}{p{\paraWidth}|}{} \\\hline
\end{tabular*}

\vspace{0.5cm}\noindent \begin{tabular*}{\tableWidth}{|c|l@{\extracolsep{\fill}}r|}
\hline
\multicolumn{1}{|p{\maxVarWidth}}{center\_offset} & {\bf Scope:} restricted & REAL \\\hline
\multicolumn{3}{|p{\descWidth}|}{{\bf Description:}   {\em offset b=0 to position (x,y,z)}} \\
\hline{\bf Range} & &  {\bf Default:} 0.0 \\\multicolumn{1}{|p{\maxVarWidth}|}{\centering (*:*)} & \multicolumn{2}{p{\paraWidth}|}{} \\\hline
\end{tabular*}

\vspace{0.5cm}\noindent \begin{tabular*}{\tableWidth}{|c|l@{\extracolsep{\fill}}r|}
\hline
\multicolumn{1}{|p{\maxVarWidth}}{do\_initial\_debug\_output} & {\bf Scope:} restricted & BOOLEAN \\\hline
\multicolumn{3}{|p{\descWidth}|}{{\bf Description:}   {\em Output debug information about initial guess}} \\
\hline & & {\bf Default:} no \\\hline
\end{tabular*}

\vspace{0.5cm}\noindent \begin{tabular*}{\tableWidth}{|c|l@{\extracolsep{\fill}}r|}
\hline
\multicolumn{1}{|p{\maxVarWidth}}{do\_residuum\_debug\_output} & {\bf Scope:} restricted & BOOLEAN \\\hline
\multicolumn{3}{|p{\descWidth}|}{{\bf Description:}   {\em Output debug information about the residuum}} \\
\hline & & {\bf Default:} no \\\hline
\end{tabular*}

\vspace{0.5cm}\noindent \begin{tabular*}{\tableWidth}{|c|l@{\extracolsep{\fill}}r|}
\hline
\multicolumn{1}{|p{\maxVarWidth}}{give\_bare\_mass} & {\bf Scope:} restricted & BOOLEAN \\\hline
\multicolumn{3}{|p{\descWidth}|}{{\bf Description:}   {\em User provides bare masses rather than target ADM masses}} \\
\hline & & {\bf Default:} yes \\\hline
\end{tabular*}

\vspace{0.5cm}\noindent \begin{tabular*}{\tableWidth}{|c|l@{\extracolsep{\fill}}r|}
\hline
\multicolumn{1}{|p{\maxVarWidth}}{grid\_setup\_method} & {\bf Scope:} restricted & KEYWORD \\\hline
\multicolumn{3}{|p{\descWidth}|}{{\bf Description:}   {\em How to fill the 3D grid from the spectral grid}} \\
\hline{\bf Range} & &  {\bf Default:} Taylor expansion \\\multicolumn{1}{|p{\maxVarWidth}|}{\centering Taylor expansion} & \multicolumn{2}{p{\paraWidth}|}{use a Taylor expansion about the nearest collocation point (fast, but might be inaccurate)} \\\multicolumn{1}{|p{\maxVarWidth}|}{\centering evaluation} & \multicolumn{2}{p{\paraWidth}|}{evaluate using all spectral coefficients (slow)} \\\hline
\end{tabular*}

\vspace{0.5cm}\noindent \begin{tabular*}{\tableWidth}{|c|l@{\extracolsep{\fill}}r|}
\hline
\multicolumn{1}{|p{\maxVarWidth}}{initial\_lapse\_psi\_exponent} & {\bf Scope:} restricted & REAL \\\hline
\multicolumn{3}{|p{\descWidth}|}{{\bf Description:}   {\em Exponent n for psi\^-n initial lapse profile}} \\
\hline{\bf Range} & &  {\bf Default:} -2.0 \\\multicolumn{1}{|p{\maxVarWidth}|}{\centering (*:*)} & \multicolumn{2}{p{\paraWidth}|}{Should be negative} \\\hline
\end{tabular*}

\vspace{0.5cm}\noindent \begin{tabular*}{\tableWidth}{|c|l@{\extracolsep{\fill}}r|}
\hline
\multicolumn{1}{|p{\maxVarWidth}}{keep\_u\_around} & {\bf Scope:} restricted & BOOLEAN \\\hline
\multicolumn{3}{|p{\descWidth}|}{{\bf Description:}   {\em Keep the variable u around after solving}} \\
\hline & & {\bf Default:} no \\\hline
\end{tabular*}

\vspace{0.5cm}\noindent \begin{tabular*}{\tableWidth}{|c|l@{\extracolsep{\fill}}r|}
\hline
\multicolumn{1}{|p{\maxVarWidth}}{multiply\_old\_lapse} & {\bf Scope:} restricted & BOOLEAN \\\hline
\multicolumn{3}{|p{\descWidth}|}{{\bf Description:}   {\em Multiply the old lapse with the new one}} \\
\hline & & {\bf Default:} no \\\hline
\end{tabular*}

\vspace{0.5cm}\noindent \begin{tabular*}{\tableWidth}{|c|l@{\extracolsep{\fill}}r|}
\hline
\multicolumn{1}{|p{\maxVarWidth}}{newton\_maxit} & {\bf Scope:} restricted & INT \\\hline
\multicolumn{3}{|p{\descWidth}|}{{\bf Description:}   {\em Maximum number of Newton iterations}} \\
\hline{\bf Range} & &  {\bf Default:} 5 \\\multicolumn{1}{|p{\maxVarWidth}|}{\centering 0:*} & \multicolumn{2}{p{\paraWidth}|}{} \\\hline
\end{tabular*}

\vspace{0.5cm}\noindent \begin{tabular*}{\tableWidth}{|c|l@{\extracolsep{\fill}}r|}
\hline
\multicolumn{1}{|p{\maxVarWidth}}{newton\_tol} & {\bf Scope:} restricted & REAL \\\hline
\multicolumn{3}{|p{\descWidth}|}{{\bf Description:}   {\em Tolerance for Newton solver}} \\
\hline{\bf Range} & &  {\bf Default:} 1.0e-10 \\\multicolumn{1}{|p{\maxVarWidth}|}{\centering (0:*)} & \multicolumn{2}{p{\paraWidth}|}{} \\\hline
\end{tabular*}

\vspace{0.5cm}\noindent \begin{tabular*}{\tableWidth}{|c|l@{\extracolsep{\fill}}r|}
\hline
\multicolumn{1}{|p{\maxVarWidth}}{npoints\_a} & {\bf Scope:} restricted & INT \\\hline
\multicolumn{3}{|p{\descWidth}|}{{\bf Description:}   {\em Number of coefficients in the compactified radial direction}} \\
\hline{\bf Range} & &  {\bf Default:} 30 \\\multicolumn{1}{|p{\maxVarWidth}|}{\centering 4:*} & \multicolumn{2}{p{\paraWidth}|}{} \\\hline
\end{tabular*}

\vspace{0.5cm}\noindent \begin{tabular*}{\tableWidth}{|c|l@{\extracolsep{\fill}}r|}
\hline
\multicolumn{1}{|p{\maxVarWidth}}{npoints\_b} & {\bf Scope:} restricted & INT \\\hline
\multicolumn{3}{|p{\descWidth}|}{{\bf Description:}   {\em Number of coefficients in the angular direction}} \\
\hline{\bf Range} & &  {\bf Default:} 30 \\\multicolumn{1}{|p{\maxVarWidth}|}{\centering 4:*} & \multicolumn{2}{p{\paraWidth}|}{} \\\hline
\end{tabular*}

\vspace{0.5cm}\noindent \begin{tabular*}{\tableWidth}{|c|l@{\extracolsep{\fill}}r|}
\hline
\multicolumn{1}{|p{\maxVarWidth}}{npoints\_phi} & {\bf Scope:} restricted & INT \\\hline
\multicolumn{3}{|p{\descWidth}|}{{\bf Description:}   {\em Number of coefficients in the phi direction}} \\
\hline{\bf Range} & &  {\bf Default:} 16 \\\multicolumn{1}{|p{\maxVarWidth}|}{\centering 4:*:2} & \multicolumn{2}{p{\paraWidth}|}{} \\\hline
\end{tabular*}

\vspace{0.5cm}\noindent \begin{tabular*}{\tableWidth}{|c|l@{\extracolsep{\fill}}r|}
\hline
\multicolumn{1}{|p{\maxVarWidth}}{par\_b} & {\bf Scope:} restricted & REAL \\\hline
\multicolumn{3}{|p{\descWidth}|}{{\bf Description:}   {\em x coordinate of the m+ puncture}} \\
\hline{\bf Range} & &  {\bf Default:} 1.0 \\\multicolumn{1}{|p{\maxVarWidth}|}{\centering (0.0:*)} & \multicolumn{2}{p{\paraWidth}|}{} \\\hline
\end{tabular*}

\vspace{0.5cm}\noindent \begin{tabular*}{\tableWidth}{|c|l@{\extracolsep{\fill}}r|}
\hline
\multicolumn{1}{|p{\maxVarWidth}}{par\_m\_minus} & {\bf Scope:} restricted & REAL \\\hline
\multicolumn{3}{|p{\descWidth}|}{{\bf Description:}   {\em mass of the m- puncture}} \\
\hline{\bf Range} & &  {\bf Default:} 1.0 \\\multicolumn{1}{|p{\maxVarWidth}|}{\centering 0.0:*)} & \multicolumn{2}{p{\paraWidth}|}{} \\\hline
\end{tabular*}

\vspace{0.5cm}\noindent \begin{tabular*}{\tableWidth}{|c|l@{\extracolsep{\fill}}r|}
\hline
\multicolumn{1}{|p{\maxVarWidth}}{par\_m\_plus} & {\bf Scope:} restricted & REAL \\\hline
\multicolumn{3}{|p{\descWidth}|}{{\bf Description:}   {\em mass of the m+ puncture}} \\
\hline{\bf Range} & &  {\bf Default:} 1.0 \\\multicolumn{1}{|p{\maxVarWidth}|}{\centering 0.0:*)} & \multicolumn{2}{p{\paraWidth}|}{} \\\hline
\end{tabular*}

\vspace{0.5cm}\noindent \begin{tabular*}{\tableWidth}{|c|l@{\extracolsep{\fill}}r|}
\hline
\multicolumn{1}{|p{\maxVarWidth}}{par\_p\_minus} & {\bf Scope:} restricted & REAL \\\hline
\multicolumn{3}{|p{\descWidth}|}{{\bf Description:}   {\em momentum of the m- puncture}} \\
\hline{\bf Range} & &  {\bf Default:} 0.0 \\\multicolumn{1}{|p{\maxVarWidth}|}{\centering (*:*)} & \multicolumn{2}{p{\paraWidth}|}{} \\\hline
\end{tabular*}

\vspace{0.5cm}\noindent \begin{tabular*}{\tableWidth}{|c|l@{\extracolsep{\fill}}r|}
\hline
\multicolumn{1}{|p{\maxVarWidth}}{par\_p\_plus} & {\bf Scope:} restricted & REAL \\\hline
\multicolumn{3}{|p{\descWidth}|}{{\bf Description:}   {\em momentum of the m+ puncture}} \\
\hline{\bf Range} & &  {\bf Default:} 0.0 \\\multicolumn{1}{|p{\maxVarWidth}|}{\centering (*:*)} & \multicolumn{2}{p{\paraWidth}|}{} \\\hline
\end{tabular*}

\vspace{0.5cm}\noindent \begin{tabular*}{\tableWidth}{|c|l@{\extracolsep{\fill}}r|}
\hline
\multicolumn{1}{|p{\maxVarWidth}}{par\_s\_minus} & {\bf Scope:} restricted & REAL \\\hline
\multicolumn{3}{|p{\descWidth}|}{{\bf Description:}   {\em spin of the m- puncture}} \\
\hline{\bf Range} & &  {\bf Default:} 0.0 \\\multicolumn{1}{|p{\maxVarWidth}|}{\centering (*:*)} & \multicolumn{2}{p{\paraWidth}|}{} \\\hline
\end{tabular*}

\vspace{0.5cm}\noindent \begin{tabular*}{\tableWidth}{|c|l@{\extracolsep{\fill}}r|}
\hline
\multicolumn{1}{|p{\maxVarWidth}}{par\_s\_plus} & {\bf Scope:} restricted & REAL \\\hline
\multicolumn{3}{|p{\descWidth}|}{{\bf Description:}   {\em spin of the m+ puncture}} \\
\hline{\bf Range} & &  {\bf Default:} 0.0 \\\multicolumn{1}{|p{\maxVarWidth}|}{\centering (*:*)} & \multicolumn{2}{p{\paraWidth}|}{} \\\hline
\end{tabular*}

\vspace{0.5cm}\noindent \begin{tabular*}{\tableWidth}{|c|l@{\extracolsep{\fill}}r|}
\hline
\multicolumn{1}{|p{\maxVarWidth}}{rescale\_sources} & {\bf Scope:} restricted & BOOLEAN \\\hline
\multicolumn{3}{|p{\descWidth}|}{{\bf Description:}   {\em If sources are used - rescale them after solving?}} \\
\hline & & {\bf Default:} yes \\\hline
\end{tabular*}

\vspace{0.5cm}\noindent \begin{tabular*}{\tableWidth}{|c|l@{\extracolsep{\fill}}r|}
\hline
\multicolumn{1}{|p{\maxVarWidth}}{schedule\_in\_admbase\_initialdata} & {\bf Scope:} restricted & BOOLEAN \\\hline
\multicolumn{3}{|p{\descWidth}|}{{\bf Description:}   {\em Schedule in (instead of after) ADMBase\_InitialData}} \\
\hline & & {\bf Default:} yes \\\hline
\end{tabular*}

\vspace{0.5cm}\noindent \begin{tabular*}{\tableWidth}{|c|l@{\extracolsep{\fill}}r|}
\hline
\multicolumn{1}{|p{\maxVarWidth}}{solve\_momentum\_constraint} & {\bf Scope:} restricted & BOOLEAN \\\hline
\multicolumn{3}{|p{\descWidth}|}{{\bf Description:}   {\em Solve for momentum constraint?}} \\
\hline & & {\bf Default:} no \\\hline
\end{tabular*}

\vspace{0.5cm}\noindent \begin{tabular*}{\tableWidth}{|c|l@{\extracolsep{\fill}}r|}
\hline
\multicolumn{1}{|p{\maxVarWidth}}{swap\_xz} & {\bf Scope:} restricted & BOOLEAN \\\hline
\multicolumn{3}{|p{\descWidth}|}{{\bf Description:}   {\em Swap x and z coordinates when interpolating, so that the black holes are separated in the z direction}} \\
\hline & & {\bf Default:} no \\\hline
\end{tabular*}

\vspace{0.5cm}\noindent \begin{tabular*}{\tableWidth}{|c|l@{\extracolsep{\fill}}r|}
\hline
\multicolumn{1}{|p{\maxVarWidth}}{target\_m\_minus} & {\bf Scope:} restricted & REAL \\\hline
\multicolumn{3}{|p{\descWidth}|}{{\bf Description:}   {\em target ADM mass for m-}} \\
\hline{\bf Range} & &  {\bf Default:} 0.5 \\\multicolumn{1}{|p{\maxVarWidth}|}{\centering 0.0:*)} & \multicolumn{2}{p{\paraWidth}|}{} \\\hline
\end{tabular*}

\vspace{0.5cm}\noindent \begin{tabular*}{\tableWidth}{|c|l@{\extracolsep{\fill}}r|}
\hline
\multicolumn{1}{|p{\maxVarWidth}}{target\_m\_plus} & {\bf Scope:} restricted & REAL \\\hline
\multicolumn{3}{|p{\descWidth}|}{{\bf Description:}   {\em target ADM mass for m+}} \\
\hline{\bf Range} & &  {\bf Default:} 0.5 \\\multicolumn{1}{|p{\maxVarWidth}|}{\centering 0.0:*)} & \multicolumn{2}{p{\paraWidth}|}{} \\\hline
\end{tabular*}

\vspace{0.5cm}\noindent \begin{tabular*}{\tableWidth}{|c|l@{\extracolsep{\fill}}r|}
\hline
\multicolumn{1}{|p{\maxVarWidth}}{tp\_epsilon} & {\bf Scope:} restricted & REAL \\\hline
\multicolumn{3}{|p{\descWidth}|}{{\bf Description:}   {\em A small number to smooth out singularities at the puncture locations}} \\
\hline{\bf Range} & &  {\bf Default:} 0.0 \\\multicolumn{1}{|p{\maxVarWidth}|}{\centering 0:*} & \multicolumn{2}{p{\paraWidth}|}{} \\\hline
\end{tabular*}

\vspace{0.5cm}\noindent \begin{tabular*}{\tableWidth}{|c|l@{\extracolsep{\fill}}r|}
\hline
\multicolumn{1}{|p{\maxVarWidth}}{tp\_extend\_radius} & {\bf Scope:} restricted & REAL \\\hline
\multicolumn{3}{|p{\descWidth}|}{{\bf Description:}   {\em Radius of an extended spacetime instead of the puncture}} \\
\hline{\bf Range} & &  {\bf Default:} 0.0 \\\multicolumn{1}{|p{\maxVarWidth}|}{\centering 0:*} & \multicolumn{2}{p{\paraWidth}|}{anything positive, should be smaller than the horizon} \\\hline
\end{tabular*}

\vspace{0.5cm}\noindent \begin{tabular*}{\tableWidth}{|c|l@{\extracolsep{\fill}}r|}
\hline
\multicolumn{1}{|p{\maxVarWidth}}{tp\_tiny} & {\bf Scope:} restricted & REAL \\\hline
\multicolumn{3}{|p{\descWidth}|}{{\bf Description:}   {\em Tiny number to avoid nans near or at the puncture locations}} \\
\hline{\bf Range} & &  {\bf Default:} 0.0 \\\multicolumn{1}{|p{\maxVarWidth}|}{\centering 0:*} & \multicolumn{2}{p{\paraWidth}|}{anything positive, usually very small} \\\hline
\end{tabular*}

\vspace{0.5cm}\noindent \begin{tabular*}{\tableWidth}{|c|l@{\extracolsep{\fill}}r|}
\hline
\multicolumn{1}{|p{\maxVarWidth}}{use\_external\_initial\_guess} & {\bf Scope:} restricted & BOOLEAN \\\hline
\multicolumn{3}{|p{\descWidth}|}{{\bf Description:}   {\em Set initial guess by external function?}} \\
\hline & & {\bf Default:} no \\\hline
\end{tabular*}

\vspace{0.5cm}\noindent \begin{tabular*}{\tableWidth}{|c|l@{\extracolsep{\fill}}r|}
\hline
\multicolumn{1}{|p{\maxVarWidth}}{use\_sources} & {\bf Scope:} restricted & BOOLEAN \\\hline
\multicolumn{3}{|p{\descWidth}|}{{\bf Description:}   {\em Use sources?}} \\
\hline & & {\bf Default:} no \\\hline
\end{tabular*}

\vspace{0.5cm}\noindent \begin{tabular*}{\tableWidth}{|c|l@{\extracolsep{\fill}}r|}
\hline
\multicolumn{1}{|p{\maxVarWidth}}{verbose} & {\bf Scope:} restricted & BOOLEAN \\\hline
\multicolumn{3}{|p{\descWidth}|}{{\bf Description:}   {\em Print screen output while solving}} \\
\hline & & {\bf Default:} no \\\hline
\end{tabular*}

\vspace{0.5cm}\noindent \begin{tabular*}{\tableWidth}{|c|l@{\extracolsep{\fill}}r|}
\hline
\multicolumn{1}{|p{\maxVarWidth}}{conformal\_storage} & {\bf Scope:} shared from STATICCONFORMAL & KEYWORD \\\hline
\end{tabular*}

\vspace{0.5cm}\parskip = 10pt 

\section{Interfaces} 


\parskip = 0pt

\vspace{3mm} \subsection*{General}

\noindent {\bf Implements}: 

twopunctures
\vspace{2mm}

\noindent {\bf Inherits}: 

admbase

staticconformal

grid
\vspace{2mm}
\subsection*{Grid Variables}
\vspace{5mm}\subsubsection{PRIVATE GROUPS}

\vspace{5mm}

\begin{tabular*}{150mm}{|c|c@{\extracolsep{\fill}}|rl|} \hline 
~ {\bf Group Names} ~ & ~ {\bf Variable Names} ~  &{\bf Details} ~ & ~\\ 
\hline 
puncture\_u & puncture\_u & compact & 0 \\ 
 &  & dimensions & 3 \\ 
 &  & distribution & DEFAULT \\ 
 &  & group type & GF \\ 
 &  & tags & prolongation="none" \\ 
 &  & timelevels & 1 \\ 
 &  & variable type & REAL \\ 
\hline 
energy &  & compact & 0 \\ 
 & E & description & ADM energy of the Bowen-York spacetime \\ 
 &  & dimensions & 0 \\ 
 &  & distribution & CONSTANT \\ 
 &  & group type & SCALAR \\ 
 &  & timelevels & 1 \\ 
 &  & variable type & REAL \\ 
\hline 
angular\_momentum &  & compact & 0 \\ 
 & J1 & description & Angular momentum of the Bowen-York spacetime \\ 
 & J2 & dimensions & 0 \\ 
 & J3 & distribution & CONSTANT \\ 
 &  & group type & SCALAR \\ 
 &  & timelevels & 1 \\ 
 &  & variable type & REAL \\ 
\hline 
\end{tabular*} 


\vspace{5mm}\subsubsection{PUBLIC GROUPS}

\vspace{5mm}

\begin{tabular*}{150mm}{|c|c@{\extracolsep{\fill}}|rl|} \hline 
~ {\bf Group Names} ~ & ~ {\bf Variable Names} ~  &{\bf Details} ~ & ~\\ 
\hline 
bare\_mass &  & compact & 0 \\ 
 & mp & description & Bare masses of the punctures \\ 
 & mm & dimensions & 0 \\ 
 &  & distribution & CONSTANT \\ 
 &  & group type & SCALAR \\ 
 &  & timelevels & 1 \\ 
 &  & variable type & REAL \\ 
\hline 
puncture\_adm\_mass &  & compact & 0 \\ 
 & mp\_adm & description & ADM masses of the punctures (measured at the other spatial infinities) \\ 
 & mm\_adm & dimensions & 0 \\ 
 &  & distribution & CONSTANT \\ 
 &  & group type & SCALAR \\ 
 &  & timelevels & 1 \\ 
 &  & variable type & REAL \\ 
\hline 
\end{tabular*} 



\vspace{5mm}

\noindent {\bf Adds header}: 



TwoPunctures.h

TP\_utilities.h

TP\_params.h

TP\_Guts.h
\vspace{2mm}\parskip = 10pt 

\section{Schedule} 


\parskip = 0pt


\noindent This section lists all the variables which are assigned storage by thorn EinsteinInitialData/TwoPunctures.  Storage can either last for the duration of the run ({\bf Always} means that if this thorn is activated storage will be assigned, {\bf Conditional} means that if this thorn is activated storage will be assigned for the duration of the run if some condition is met), or can be turned on for the duration of a schedule function.


\subsection*{Storage}

\hspace{5mm}

 \begin{tabular*}{160mm}{ll} 
~& {\bf Conditional:} \\ 
~ &  energy angular\_momentum puncture\_adm\_mass bare\_mass\\ 
~ &  puncture\_u\\ 
~ & ~\\ 
\end{tabular*} 


\subsection*{Scheduled Functions}
\vspace{5mm}

\noindent {\bf CCTK\_PARAMCHECK}   (conditional) 

\hspace{5mm} twopunctures\_paramcheck 

\hspace{5mm}{\it check parameters and thorn needs } 


\hspace{5mm}

 \begin{tabular*}{160mm}{cll} 
~ & Language:  & c \\ 
~ & Type:  & function \\ 
\end{tabular*} 


\vspace{5mm}

\noindent {\bf ADMBase\_InitialData}   (conditional) 

\hspace{5mm} twopunctures\_group 

\hspace{5mm}{\it twopunctures initial data group } 


\hspace{5mm}

 \begin{tabular*}{160mm}{cll} 
~ & Type:  & group \\ 
\end{tabular*} 


\vspace{5mm}

\noindent {\bf CCTK\_INITIAL}   (conditional) 

\hspace{5mm} twopunctures\_group 

\hspace{5mm}{\it twopunctures initial data group } 


\hspace{5mm}

 \begin{tabular*}{160mm}{cll} 
~ & After:  & admbase\_initialdata \\ 
~& ~ &hydrobase\_initial\\ 
~ & Before:  & admbase\_postinitial \\ 
~& ~ &settmunu\\ 
~& ~ &hydrobase\_prim2coninitial\\ 
~ & Type:  & group \\ 
\end{tabular*} 


\vspace{5mm}

\noindent {\bf TwoPunctures\_Group}   (conditional) 

\hspace{5mm} twopunctures 

\hspace{5mm}{\it create puncture black hole initial data } 


\hspace{5mm}

 \begin{tabular*}{160mm}{cll} 
~ & Language:  & c \\ 
~ & Reads:  & grid::coordinates(everywhere) \\ 
~ & Storage:  & puncture\_u \\ 
~ & Type:  & function \\ 
~ & Writes:  & twopunctures::mp(everywhere) \\ 
~& ~ &twopunctures::mm(everywhere)\\ 
~& ~ &twopunctures::mp\_adm(everywhere)\\ 
~& ~ &twopunctures::mm\_adm(everywhere)\\ 
~& ~ &twopunctures::e(everywhere)\\ 
~& ~ &twopunctures::j1(everywhere)\\ 
~& ~ &twopunctures::j2(everywhere)\\ 
~& ~ &twopunctures::j3(everywhere)\\ 
~& ~ &twopunctures::puncture\_u(everywhere)\\ 
~& ~ &staticconformal::conformal\_state(everywhere)\\ 
~& ~ &staticconformal::confac\_2derivs(everywhere)\\ 
~& ~ &staticconformal::confac\_1derivs(everywhere)\\ 
~& ~ &staticconformal::confac(everywhere)\\ 
~& ~ &admbase::alp(everywhere)\\ 
~& ~ &admbase::metric(everywhere)\\ 
~& ~ &admbase::curv(everywhere)\\ 
\end{tabular*} 


\vspace{5mm}

\noindent {\bf TwoPunctures\_Group}   (conditional) 

\hspace{5mm} twopunctures\_metadata 

\hspace{5mm}{\it output twopunctures metadata } 


\hspace{5mm}

 \begin{tabular*}{160mm}{cll} 
~ & After:  & twopunctures \\ 
~ & Language:  & c \\ 
~ & Options:  & global \\ 
~ & Type:  & function \\ 
\end{tabular*} 



\vspace{5mm}\parskip = 10pt 
\end{document}
